\documentclass[11pt]{article}
\usepackage{cmap}            % copyable/searchable text: ToUnicode CMaps
\usepackage[margin=1.1in]{geometry}
\usepackage{amsmath,amssymb,amsthm}
\usepackage[T1]{fontenc}
\usepackage{lmodern}
\input{glyphtounicode}
\pdfgentounicode=1
\usepackage{xcolor}
\usepackage{hyperref}
\definecolor{linkblue}{rgb}{0.1,0.2,0.55}
\hypersetup{colorlinks=true,linkcolor=linkblue,citecolor=linkblue,urlcolor=linkblue}

\newtheorem{theorem}{Theorem}
\newtheorem{lemma}[theorem]{Lemma}
\newtheorem{corollary}[theorem]{Corollary}
\newtheorem{proposition}[theorem]{Proposition}
\theoremstyle{remark}
\newtheorem{remark}[theorem]{Remark}

\title{Sharp Continuous-Order Comparison for Modified-Bessel Ratios:
a Riccati Proof, the Optimal Real-Order Domain, and Lean Verification}
\author{Lluis Eriksson}
\date{August 2026}

\begin{document}
\maketitle

\begin{abstract}
For $x>0$ put $\rho_\alpha(x)=I_{\alpha+1}(x)/I_\alpha(x)$, where
$I_\alpha$ is the modified Bessel function of the first kind.  We study
the sharp continuous-order comparison
\[
  0<\rho_\mu(x)-\rho_\nu(x)<\frac{\nu-\mu}{x}.
\]
The lower inequality is the classical strict order decrease of the
Bessel ratio.  For the upper inequality we give a direct two-flow
Riccati proof.  If $h=\rho_\mu-\rho_\nu$, $b=(\nu-\mu)/x$, and
$h=b$ at a contact point, then
\[
 h'=-\frac{(\nu-\mu)(\mu+\nu+1)}{x^2},\qquad
 (h-b)'=-\frac{(\nu-\mu)(\mu+\nu)}{x^2}.
\]
The second identity, not the first, is the contact orientation.  It
proves the bound for $\mu+\nu>0$.  The degenerate boundary
$\mu+\nu=0$ follows from the connection formula between $I_{-a}$,
$I_a$, and $K_a$ together with their Wronskian.  Conversely the
large-$x$ expansion gives counterexamples whenever $\mu+\nu<0$.
Thus, on the natural positivity range $-1<\mu<\nu$, the upper bound
holds for every $x>0$ if and only if $\mu+\nu\ge0$.

The nonnegative-order statement itself is known: it is equivalent to
Lemma~10 of Freitas--Laugesen.  Their proof already contains a
first-contact argument, so no novelty is claimed for that mechanism.
The contributions here are the explicit sharp ratio formulation, the
optimal extension/obstruction across negative orders, the correction
of an earlier unit-step paper statement, and a Lean~4/Mathlib
formalization of the complete nonnegative-order upper bound.  The Lean
artifact has no \texttt{sorry}; its axiom oracle returns exactly
\texttt{propext}, \texttt{Classical.choice}, and \texttt{Quot.sound}.
We retain the resulting Feynman--Hellmann application: every character
sector gap in two-dimensional Wilson $U(1)$ and $SU(2)$ lattice gauge
theory decreases strictly with the bare coupling.
\end{abstract}

\section{Introduction}

The ratios $\rho_\nu(x) = I_{\nu+1}(x)/I_\nu(x)$ of modified Bessel
functions of consecutive orders appear throughout applied analysis,
probability and statistical physics, and a substantial literature is
devoted to bounding them by algebraic functions of $\nu$ and $x$
\cite{Amos, Nasell, Segura2011, HornikGrun, RuizSegura, Segura2023}.
This note concerns the behaviour under arbitrary real order shifts of the logarithmic
derivative
\[
\frac{\partial}{\partial x} \log I_\nu(x) \;=\; \frac{I_\nu'(x)}{I_\nu(x)}.
\]

\begin{theorem}[sharp continuous-order comparison]
\label{thm:main}
Let $x>0$ and $0\le \mu<\nu$.  Then
\begin{equation}
\label{eq:quantitative}
0<\rho_\mu(x)-\rho_\nu(x)<\frac{\nu-\mu}{x}.
\end{equation}
Equivalently,
\begin{equation}
\label{eq:continuouslog}
\frac{I_\mu'(x)}{I_\mu(x)}<\frac{I_\nu'(x)}{I_\nu(x)}.
\end{equation}
The left inequality is classical; the proof below addresses the sharp
right inequality directly.
\end{theorem}

\begin{theorem}[optimal real-order domain]
\label{thm:domain}
Let $-1<\mu<\nu$.  The lower inequality in
\eqref{eq:quantitative} holds for every $x>0$.  The upper inequality
holds for every $x>0$ if and only if
\[
  \mu+\nu\ge0.
\]
If $\mu+\nu<0$, then the upper inequality fails for all sufficiently
large $x$.
\end{theorem}

\begin{remark}[priority and corrected scope]
\label{rem:scope}
The nonnegative-order conclusion \eqref{eq:continuouslog} is Lemma~10
of Freitas and Laugesen \cite{FreitasLaugesen}.  A previous version of
this note stated that fractional steps did not follow from its method
and described their proof as a Bessel-zero argument.  Both statements
were wrong.  Their published proof defines the logarithmic-derivative
difference, uses its positive small-$x$ sign, and excludes a first zero
directly from the modified-Bessel differential equation.  The
two-flow calculation below is the same contact geometry written for
$\rho_\mu$ and $\rho_\nu$.  We therefore claim neither the
nonnegative-order theorem nor the first-contact mechanism as new.
Segura's Riccati/nullcline analysis \cite{Segura2021} is also close in
method, although its stated monotonicity results vary $x$ at fixed
order rather than compare two arbitrary orders.  The genuinely new
candidate isolated here is the optimal negative-order boundary in
Theorem~\ref{thm:domain}, together with the formal verification.
\end{remark}

The purpose of this corrected note is threefold.

\emph{First}, we expose the entire fractional-order comparison as one
Riccati contact identity.  This avoids Bessel zeros and high-order
uniform asymptotics in the proof region $\mu+\nu>0$.

\emph{Second}, we resolve the contact-degenerate line $\mu+\nu=0$ and
show by asymptotics that crossing it reverses the large-$x$ sign.  This
produces an exact global domain rather than a sufficient condition.

\emph{Third}, we record a consequence in two-dimensional lattice gauge
theory (we are not aware of an explicit statement in the literature,
though it may well be folklore): by the
Feynman--Hellmann identity --- exact here, because the 2D transfer
operator is diagonal in a $\beta$-independent character basis --- the
continuous-order comparison of Theorem~\ref{thm:main} is precisely the strict decrease in
the bare coupling $\beta$ of every mass gap between character sectors,
for the $U(1)$ and $SU(2)$ Wilson actions simultaneously
(Section~\ref{sec:ym}). We emphasize the modest scope: this is a
statement about the exactly soluble two-dimensional theory; nothing
here is claimed about four-dimensional gauge theory.

Section~\ref{sec:lean} presents the machine-checked Lean~4/Mathlib
verification of the full nonnegative-order upper theorem (no
\texttt{sorry}; standard axioms only).
Section~\ref{sec:numerics} reports an independent numerical audit.

\section{The two-flow Riccati proof and the optimal domain}
\label{sec:proof}

Write $\delta=\nu-\mu>0$ and
\[
 h(x)=\rho_\mu(x)-\rho_\nu(x),\qquad b(x)=\frac{\delta}{x},
 \qquad G(x)=h(x)-b(x).
\]
The derivative and recurrence identities imply the Riccati equation
\begin{equation}
\label{eq:riccati}
 \rho_\alpha'=1-\rho_\alpha^2-\frac{2\alpha+1}{x}\rho_\alpha.
\end{equation}
Subtracting the equations at $\mu$ and $\nu$ gives
\begin{equation}
\label{eq:twoflow}
 h'=-(\rho_\mu+\rho_\nu)h-\frac{2\mu+1}{x}h
       +\frac{2\delta}{x}\rho_\nu.
\end{equation}
At a point where $h=b=\delta/x$, substitute
$\rho_\mu=\rho_\nu+\delta/x$ into \eqref{eq:twoflow}.  Direct
cancellation yields
\begin{equation}
\label{eq:touch}
 h'=-\frac{\delta(\mu+\nu+1)}{x^2},\qquad
 G'=h'-b'=-\frac{\delta(\mu+\nu)}{x^2}.
\end{equation}
This distinction is essential: the first displayed derivative is the
derivative of the ratio difference, whereas first-contact orientation
is controlled by the second.

For $\alpha>-1$, the defining series gives
\begin{equation}
\label{eq:smallx}
 \rho_\alpha(x)=\frac{x}{2(\alpha+1)}
 \left(1-\frac{x^2}{4(\alpha+1)(\alpha+2)}+O(x^4)\right).
\end{equation}
Consequently $h(x)>0$ for small $x$, while $G(x)\to-\infty$ as
$x\downarrow0$.  If $\mu+\nu>0$ and $G$ had a first zero $c$, then
$G<0$ immediately to its left but \eqref{eq:touch} would give
$G'(c)<0$.  Differentiability would instead force $G>G(c)=0$
immediately to the left, a contradiction.  Hence $G<0$ for all
$x>0$.  This proves the upper half of Theorem~\ref{thm:main} and the
sufficiency assertion of Theorem~\ref{thm:domain} away from the
boundary.  Notice also that at any putative zero of $h$,
\eqref{eq:twoflow} gives
\[
 h'=\frac{2\delta}{x}\rho_\nu>0.
\]
Together with the small-$x$ sign this is the analogous first-contact
orientation for the lower inequality.  More generally, strict order
decrease of $\alpha\mapsto\rho_\alpha(x)$ on $(-1,\infty)$ is a
classical consequence of order log-concavity; see
\cite{BariczRatio}.  We use that standard result for the full lower
statement.

It remains to handle the degenerate line $\mu+\nu=0$.  Write
$\mu=-a$, $\nu=a$, where $0<a<1$.  The connection formula
\cite[10.27.4]{DLMF}
\begin{equation}
\label{eq:connection}
 I_{-a}(x)=I_a(x)+\frac{2}{\pi}\sin(\pi a)K_a(x)
\end{equation}
has a strictly positive coefficient of $K_a$.  The Wronskian
\cite[10.28.2]{DLMF} gives
\begin{equation}
\label{eq:wronskian}
 \left(\frac{K_a}{I_a}\right)'
 =\frac{K_a'I_a-K_aI_a'}{I_a^2}
 =-\frac{1}{xI_a(x)^2}<0.
\end{equation}
Thus $I_a/I_{-a}=\bigl(1+c_aK_a/I_a\bigr)^{-1}$ is strictly
increasing.  Equivalently,
$(\log I_a)'-(\log I_{-a})'>0$, which by
$(\log I_\alpha)'=\rho_\alpha+\alpha/x$ is precisely
$\rho_{-a}-\rho_a<2a/x$.  This proves strictness on the boundary.

Finally, the standard large-$x$ expansion is
\begin{equation}
\label{eq:largex}
 \rho_\alpha(x)=1-\frac{\alpha+\tfrac12}{x}
 +\frac{4\alpha^2-1}{8x^2}+O(x^{-3}).
\end{equation}
Hence
\begin{equation}
\label{eq:obstruction}
 h(x)-\frac{\delta}{x}
 =-\frac{\delta(\mu+\nu)}{2x^2}+O(x^{-3}).
\end{equation}
When $\mu+\nu<0$, the right-hand side is positive for all sufficiently
large $x$, proving necessity and completing Theorem~\ref{thm:domain}.

For a concrete audit witness, at $100$ decimal digits,
\[
 (\mu,\nu,x)=(-0.8,-0.4,10),\qquad
 \rho_\mu-\rho_\nu=0.042682961675102858\ldots>0.04=\frac{\nu-\mu}{x}.
\]

\section{The unit-step Amos calibration}
\label{sec:unitstep}

Throughout, $x>0$, $\nu \ge 0$, $I_\nu$ is the modified Bessel function
of the first kind, all values $I_\nu(x)$ are strictly positive, and
$\rho_\nu(x) = I_{\nu+1}(x)/I_\nu(x) \in (0,1)$. We use two classical
facts, valid for all real orders (see \cite[10.29.1--10.29.2]{DLMF}):
\begin{align}
I_{\nu-1}(x) - I_{\nu+1}(x) &= \frac{2\nu}{x}\, I_\nu(x),
\label{eq:recurrence}\\
I_\nu'(x) &= \tfrac12\bigl(I_{\nu-1}(x) + I_{\nu+1}(x)\bigr).
\label{eq:derivative}
\end{align}

Specializing Theorem~\ref{thm:main} to the unit shift gives
\begin{equation}
\label{eq:unitstep}
\frac{I_{\nu+1}'(x)}{I_{\nu+1}(x)}-
\frac{I_\nu'(x)}{I_\nu(x)}
=\frac1x-\bigl(\rho_\nu(x)-\rho_{\nu+1}(x)\bigr)>0.
\end{equation}
The older four-step proof of this specialization is retained because
it records the exact equivalence with the Amos bound.

\subsection*{Step 1: a closed form for the logarithmic derivative}

Dividing \eqref{eq:derivative} by $I_\nu$ and using
\eqref{eq:recurrence} in the form $I_{\nu-1}/I_\nu = \rho_\nu + 2\nu/x$:
\begin{equation}
\label{eq:closedform}
\frac{I_\nu'(x)}{I_\nu(x)}
= \frac12\left(\frac{I_{\nu-1}}{I_\nu} + \frac{I_{\nu+1}}{I_\nu}\right)
= \frac12\left(\rho_\nu + \frac{2\nu}{x} + \rho_\nu\right)
= \rho_\nu(x) + \frac{\nu}{x}.
\end{equation}

\subsection*{Step 2: reduction to a ratio inequality}

By \eqref{eq:closedform}, for the unit step $\nu \to \nu+1$,
\[
\frac{I_{\nu+1}'}{I_{\nu+1}} - \frac{I_\nu'}{I_\nu}
= \rho_{\nu+1} - \rho_\nu + \frac{1}{x},
\]
which is \eqref{eq:unitstep}. Writing the recurrence
\eqref{eq:recurrence} at order $\nu+1$ as
$\rho_{\nu+1} = 1/\rho_\nu - 2(\nu+1)/x$, the inequality
$\rho_\nu - \rho_{\nu+1} < 1/x$ is equivalent to
\begin{equation}
\label{eq:target}
\frac{1}{\rho_\nu(x)} - \rho_\nu(x) \;>\; \frac{2\nu+1}{x}.
\end{equation}

\subsection*{Step 3: the Amos-type bound}

We use the classical upper bound (Amos \cite{Amos}; proved best of its
form and valid for all $\nu \ge 0$ by Segura
\cite[Thm.~2]{RuizSegura}, after the index shift $\nu \mapsto \nu+1$):
\begin{equation}
\label{eq:amos}
\rho_\nu(x) \;<\; U_\nu(x) := \frac{x}{a + s},
\qquad a = \nu + \tfrac12,\quad s = \sqrt{a^2 + x^2}.
\end{equation}

\subsection*{Step 4: the exact calibration}

\begin{lemma}[calibration]
\label{lem:calibration}
Let $a, x > 0$ and $s = \sqrt{a^2+x^2}$, $U = x/(a+s)$. Then
\[
\frac{1}{U} - U = \frac{2a}{x},
\]
and for every $t$ with $0 < t < U$,
\[
\frac{1}{t} - t \;>\; \frac{2a}{x}.
\]
\end{lemma}

\begin{proof}
Since $(s-a)(s+a) = s^2 - a^2 = x^2$, we have
$U = x/(a+s) = (s-a)/x$, hence
\[
\frac1U - U = \frac{a+s}{x} - \frac{s-a}{x} = \frac{2a}{x}.
\]
The map $t \mapsto 1/t - t$ is strictly decreasing on $(0,\infty)$, so
$t < U$ gives $1/t - t > 1/U - U = 2a/x$.
\end{proof}

Applying Lemma~\ref{lem:calibration} with $t = \rho_\nu(x) < U_\nu(x)$
and $a = \nu + \frac12$ yields \eqref{eq:target} strictly, hence
\eqref{eq:unitstep} and the unit-shift instance of
\eqref{eq:quantitative}. \qed

\begin{remark}[equivalence with the Amos bound]
\label{rem:equivalence}
The chain of Step~2 consists of equivalences: since
$t \mapsto 1/t - t$ is a strictly decreasing bijection of $(0,\infty)$
onto $\mathbb{R}$, and $U_\nu$ is precisely the positive root of
$1/t - t = (2\nu+1)/x$ (Lemma~\ref{lem:calibration}), the inequality
\eqref{eq:target} is \emph{equivalent} to $\rho_\nu < U_\nu$.
Therefore: \emph{the unit-step order-monotonicity \eqref{eq:unitstep}
is equivalent to the classical Amos upper bound \eqref{eq:amos}}. Any
of the many proofs of the latter \cite{Amos, Segura2011, HornikGrun,
RuizSegura} yields the former, and conversely the continuous
monotonicity of Freitas--Laugesen \cite{FreitasLaugesen} yields, in
particular, a Bessel-zero proof of the Amos bound. The bound $U_\nu$
is also recognizable as the $b_0$-type bound of \cite{Segura2011}: the
positive root of the Riccati-derived quadratic
$t^2 + \frac{2\nu+1}{x}t - 1 = 0$.
\end{remark}

\section{Application: bare-coupling monotonicity of all 2D sector gaps}
\label{sec:ym}

Consider lattice gauge theory on a two-dimensional rectangular lattice
with Wilson action and compact structure group $G$, at bare coupling
$\beta > 0$. We fix notation carefully, separating the three objects
involved.

\emph{(a) Character coefficients.} The one-plaquette Boltzmann weight
expands as (see, e.g., \cite{Migdal, DrouffeZuber})
\[
e^{\beta\,\mathrm{Re}\,\chi_f(U)}
\;=\; \sum_r d_r\, c_r(\beta)\, \chi_r(U),
\qquad
c_r(\beta) = \frac{1}{d_r}\int_G \overline{\chi_r(U)}\,
e^{\beta\,\mathrm{Re}\,\chi_f(U)}\, dU ,
\]
the sum over irreducible representations $r$ of $G$, $d_r = \chi_r(e)$.
For the two classical cases ($n \in \mathbb{Z}_{\ge0}$,
$j \in \tfrac12\mathbb{Z}_{\ge0}$; verified independently to 30 digits
in the audit of Section~\ref{sec:numerics}):
\[
G = U(1):\quad c_n(\beta) = I_n(\beta), \qquad\quad
G = SU(2):\quad c_j(\beta) = \frac{I_{2j+1}(2\beta)}{\beta},
\]
the $SU(2)$ formula being the statement that the Haar class integral
equals $\int_G \chi_j(U)\, e^{\beta\chi_f(U)}\,dU = (2j+1)\,
I_{2j+1}(2\beta)/\beta$, i.e.\ $d_j\,c_j(\beta)$ with $d_j = 2j+1$
(all $SU(2)$ characters are real, so $\mathrm{Re}\,\chi_f = \chi_f
= 2\cos\theta$).

\emph{(b) Transfer eigenvalues.} In two dimensions, gauge fixing
reduces the theory to a product of independent plaquette weights, and
the transfer operator (one lattice row, temporal gauge) is diagonal in
the character basis $\{\chi_r\}$ --- the standard exact solubility of
2D lattice gauge theory \cite{Migdal, GrossWitten}; see
\cite{MontvayMunster} for a textbook treatment of the character
expansion --- a basis independent of $\beta$ ---
with eigenvalue in sector $r$ proportional to $c_r(\beta)$, the
proportionality constant being a sector-independent normalization.
Ratios of transfer eigenvalues between sectors are therefore exactly
$c_r(\beta)/c_{r'}(\beta)$, with every normalization convention
cancelling.

\emph{(c) Sector gaps.} For two sectors with
$c_r(\beta) > c_{r'}(\beta) > 0$ (in both models the coefficient is
strictly decreasing in the sector label, so ``heavier'' sectors have
smaller coefficients), the mass gap of sector $r'$ relative to $r$ is
\[
m_{r,r'}(\beta) \;=\; \log\frac{c_r(\beta)}{c_{r'}(\beta)} \;>\; 0,
\]
e.g.\ for $U(1)$ the fundamental string tension is
$m_{0,1} = \log(I_0/I_1)$ and the $v87$ bracket instance is
$m_{1,2} = \log(I_1/I_2)$. Because the eigenbasis is
$\beta$-independent, $\frac{d}{d\beta} m_{r,r'}$ acts on the
eigenvalues alone --- the Feynman--Hellmann identity in its exact,
diagonal form; no perturbation theory is involved.

\begin{corollary}
\label{cor:ym}
For 2D Wilson lattice gauge theory:
\begin{itemize}
\item[(i)] $G = U(1)$: for all integers $0 \le n < n'$,
\[
\frac{d}{d\beta}\, m_{n,n'}(\beta)
= \frac{d}{d\beta}\, \log\frac{I_n(\beta)}{I_{n'}(\beta)}
= \bigl(\rho_n - \rho_{n'}\bigr)(\beta) - \frac{n'-n}{\beta}
\;<\; 0 .
\]
Every sector gap --- in particular the string tension $\log(I_0/I_1)$
and the bracket instance $\log(I_1/I_2)$ --- is strictly decreasing
in $\beta$.
\item[(ii)] $G = SU(2)$: for all $0 \le j < j'$ in
$\tfrac12\mathbb{Z}$,
\[
\frac{d}{d\beta}\, m_{j,j'}(\beta)
= 2\left[\bigl(\log I_{2j+1}\bigr)'(2\beta)
- \bigl(\log I_{2j'+1}\bigr)'(2\beta)\right] \;<\; 0 ,
\]
the $1/\beta$ prefactors cancelling in the ratio
$c_j/c_{j'} = I_{2j+1}(2\beta)/I_{2j'+1}(2\beta)$, which is free of
representation-dimension factors.
\end{itemize}
Both assertions are direct instances of Theorem~\ref{thm:main}; the
orders happen to be integer-spaced, so the older unit-step proof in
Section~\ref{sec:unitstep} would also suffice.
\end{corollary}

\begin{proof}
(i) By \eqref{eq:closedform},
$\frac{d}{d\beta}\log(I_n/I_{n'}) = (\rho_n + n/\beta) - (\rho_{n'} +
n'/\beta) = (\rho_n - \rho_{n'}) - (n'-n)/\beta$, and
\eqref{eq:quantitative} gives
$\rho_n - \rho_{n'} < (n'-n)/\beta$ strictly. (ii) Identical, at
argument $x = 2\beta$ and orders $2j+1 < 2j'+1$; the chain rule
contributes the factor $2$.
\end{proof}

\begin{remark}[scope]
Corollary~\ref{cor:ym} is a statement about the exactly soluble
two-dimensional theory. The analogous monotonicity $m'(\beta) \le 0$
for the four-dimensional mass gap on a weak-to-strong coupling corridor
is an open problem (universally expected in lattice physics, covered by
no theorem); the present corollary is the exact 2D instance of that
bracket and nothing more. In the terminology of the audit-first
programme to which this note belongs, the 2D instance is
\emph{signed}; the 4D bracket remains a named hypothesis.
\end{remark}

\section{The machine-checked algebraic core}
\label{sec:lean}

\begin{sloppypar}
The continuous-order upper theorem on $0\le\mu<\nu$ is formalized in
the C4 module \texttt{FractionalOrder.lean}.  Unlike the earlier
standalone algebraic artifact, this file works on C4's actual
Gamma-series definition of $I_\nu$, invokes the proved real-order
Riccati equation and Amos bound, supplies a uniform small-$x$ seed, and
executes the $\inf$ first-contact argument.  It contains no
\texttt{sorry} and introduces no project axiom.  The checked endpoint
is:
\end{sloppypar}
\begin{verbatim}
theorem besselRatioReal_fractional_upper
    (mu nu : Real) (hmu : 0 <= mu) (hmunu : mu < nu)
    {x : Real} (hx : 0 < x) :
    besselRatioReal mu x - besselRatioReal nu x
      < (nu - mu) / x

theorem besselIReal_logDeriv_fractional_lt
    (mu nu : Real) (hmu : 0 <= mu) (hmunu : mu < nu)
    {x : Real} (hx : 0 < x) :
    deriv (fun y => Real.log (besselIReal mu y)) x
      < deriv (fun y => Real.log (besselIReal nu y)) x
\end{verbatim}
The implementation exposes three audit lemmas: the derivative of the
signed gap, its exact derivative at contact, and the small-$x$ seed.
The axiom oracle for all five declarations returns exactly:
\begin{verbatim}
[propext, Classical.choice, Quot.sound]
\end{verbatim}
Verification used Lean \texttt{4.29.0-rc6} and pinned Mathlib commit
\begin{center}\small
\texttt{07642720480157414db592fa85b626dafb71355b}.
\end{center}
On 4 August 2026
the module was built in a Colab Pro+ CPU/high-memory runtime: 8165 Lake
jobs completed, the module's \texttt{.olean} was materialized, and the
oracle file compiled.  Runtime opening and teardown were recorded; the
runtime was disconnected and deleted after 18 minutes 53 seconds.

\emph{Honest scope.}  The machine-checked theorem covers the sharp
upper inequality on the nonnegative order axis.  The classical lower
inequality and the negative-order boundary/obstruction in
Theorem~\ref{thm:domain} are proved on paper but are not represented in
C4, whose real-order Bessel interface currently assumes nonnegative
order.  Source presence is not reported as verification: the claim
above refers to the materialized Colab build and its oracle output.

\section{Numerical audit}
\label{sec:numerics}

The fractional-order statement was attacked independently at 70
decimal digits on 800 deterministic adversarial points (seed
\texttt{490060}), including tiny order separations, values near
$\mu+\nu=0$, and arguments from the small- to large-$x$ regimes.  No
failure occurred in the valid nonnegative domain; this is diagnostic,
not proof.  The same script verifies the explicit invalid-domain
witness following \eqref{eq:obstruction}.  In addition, the older
unit-step/application inequalities were checked at 50 decimal digits
(mpmath), on the grid $\beta \in \{10^{-2}, \dots, 10^4\}$ (16 values,
geometrically spread), $\nu \in \{0, 1, \dots, 29\} \cup \{50, 100,
500\}$: the closed form \eqref{eq:closedform} holds to $1.4\cdot
10^{-18}$ (central finite difference at step $10^{-20}$); the Amos
margin $U_\nu - \rho_\nu$ is strictly positive throughout, with minimum
$2.5\cdot 10^{-9}$ attained at $(\beta,\nu) = (10^4, 0)$ (the bound is
asymptotically sharp, as it must be, being extremal of its form); the
quantitative slack $1/\beta - (\rho_\nu - \rho_{\nu+1})$ is strictly
positive throughout, with minimum $1.5\cdot10^{-8}$ at $(10^4, 1)$; and
$m(\beta) = \log(I_1/I_2)$ has $m'(\beta) < 0$ at every grid point,
ranging from $-1.96$ at $\beta = \tfrac12$ to $-1.9\cdot10^{-5}$ at
$\beta = 278$. The $SU(2)$ Haar class integral
$\int_G \chi_j\, e^{\beta\chi_f}\,dU = \frac{2}{\pi}\int_0^\pi
\sin\theta\, \sin((2j+1)\theta)\, e^{2\beta\cos\theta}\, d\theta$
was verified to equal $(2j+1)\, I_{2j+1}(2\beta)/\beta$ to 30 digits
at $j \in \{0,1,2\}$, $\beta \in \{0.7, 2\}$ --- hence
$c_j = I_{2j+1}(2\beta)/\beta$ after dividing by $d_j = 2j+1$. The verification script and the
sources of this note are archived in the programme repository:
\begin{center}\small
\url{https://github.com/lluiseriksson/THE-ERIKSSON-PROGRAMME/tree/main/papers/bessel-amos-fh}
\end{center}

\subsection*{Version and audit note}
This note is part of an audit-first research programme whose records,
Lean sources and numerical suites are public at
\url{https://github.com/lluiseriksson/THE-ERIKSSON-PROGRAMME}
(machine-checked core) and the satellite repositories linked there. The author thanks the
adversarial-audit process itself.  It first located
\cite{FreitasLaugesen}, correcting a novelty claim, and later exposed a
second error in version~6: the paper said its method did not close
fractional steps and misdescribed Lemma~10 as a Bessel-zero proof.  The
contact calculation \eqref{eq:touch} corrects that statement.  This
version separates the known nonnegative-order theorem and its known
first-contact route from the candidate-new optimal negative-order
classification.  Numerical residuals and a finite mesh are never used
as proof.

\begin{thebibliography}{10}

\bibitem{Amos}
D.~E. Amos,
\emph{Computation of modified Bessel functions and their ratios},
Math. Comp. \textbf{28} (1974), 239--251.

\bibitem{Nasell}
I.~N{\aa}sell,
\emph{Rational bounds for ratios of modified Bessel functions},
SIAM J. Math. Anal. \textbf{9} (1978), 1--11.

\bibitem{Segura2011}
J.~Segura,
\emph{Bounds for ratios of modified Bessel functions and associated
Tur\'an-type inequalities},
J. Math. Anal. Appl. \textbf{374} (2011), 516--528.

\bibitem{HornikGrun}
K.~Hornik and B.~Gr\"un,
\emph{Amos-type bounds for modified Bessel function ratios},
J. Math. Anal. Appl. \textbf{408} (2013), 91--101.

\bibitem{RuizSegura}
D.~Ruiz-Antol\'in and J.~Segura,
\emph{A new type of sharp bounds for ratios of modified Bessel
functions},
J. Math. Anal. Appl. \textbf{443} (2016), 1232--1246;
arXiv:1606.02008.

\bibitem{Segura2023}
J.~Segura,
\emph{Simple bounds with best possible accuracy for ratios of modified
Bessel functions},
J. Math. Anal. Appl. \textbf{526} (2023), 127211; arXiv:2207.02713.

\bibitem{FreitasLaugesen}
P.~Freitas and R.~S. Laugesen,
\emph{From Neumann to Steklov and beyond, via Robin: the Weinberger
way},
Amer. J. Math. \textbf{143} (2021), 969--994;
doi:10.1353/ajm.2021.0024; arXiv:1810.07461; Lemma~10.

\bibitem{Segura2021}
J.~Segura,
\emph{Monotonicity properties for ratios and products of modified
Bessel functions and sharp trigonometric bounds},
Results Math. \textbf{76} (2021), article 221;
doi:10.1007/s00025-021-01531-1; arXiv:2105.02524.

\bibitem{BariczTuran}
\'A.~Baricz,
\emph{On Tur\'an type inequalities for modified Bessel functions},
Proc. Amer. Math. Soc. \textbf{141} (2013), 523--532; arXiv:1010.3346.

\bibitem{BariczRatio}
\'A.~Baricz,
\emph{Bounds for Tur\'anians of modified Bessel functions},
Expo. Math. \textbf{33} (2015), 223--251;
doi:10.1016/j.exmath.2014.07.001; arXiv:1202.4853.

\bibitem{GrossWitten}
D.~J. Gross and E.~Witten,
\emph{Possible third-order phase transition in the large-$N$ lattice
gauge theory},
Phys. Rev. D \textbf{21} (1980), 446--453.

\bibitem{MontvayMunster}
I.~Montvay and G.~M\"unster,
\emph{Quantum Fields on a Lattice},
Cambridge Monographs on Mathematical Physics, Cambridge University
Press, 1994.

\bibitem{DLMF}
F.~W.~J. Olver et al.,
\emph{NIST Digital Library of Mathematical Functions},
\S\S10.27--10.29, \url{https://dlmf.nist.gov/10}.

\bibitem{Migdal}
A.~A. Migdal,
\emph{Recursion equations in gauge field theories},
Sov. Phys. JETP \textbf{42} (1975), 413--418.

\bibitem{DrouffeZuber}
J.-M. Drouffe and J.-B. Zuber,
\emph{Strong coupling and mean field methods in lattice gauge theories},
Phys. Rep. \textbf{102} (1983), 1--119.

\end{thebibliography}

\end{document}
